% !TeX program = xelatex
%-------------------------
% Resume in Latex
% Author : Jake Gutierrez
% Based off of: https://github.com/sb2nov/resume
% License : MIT
%------------------------

\documentclass[letterpaper,11pt]{article}

\usepackage{iftex}
\ifPDFTeX
  \usepackage[T2A]{fontenc}
  \usepackage[utf8]{inputenc}
  \usepackage{charter}
\else
  \usepackage{fontspec}
  \defaultfontfeatures{Ligatures=TeX}
  \setmainfont{Libertinus Serif}
\fi
\usepackage{latexsym}
\usepackage[empty]{fullpage}
\usepackage{titlesec}
\usepackage{marvosym}
\usepackage[usenames,dvipsnames]{color}
\usepackage{verbatim}
\usepackage{enumitem}
\usepackage[hidelinks]{hyperref}
\usepackage{fancyhdr}
\usepackage[russian]{babel}
\usepackage{tabularx}
\ifPDFTeX
  \input{glyphtounicode}
\fi


%----------FONT OPTIONS----------
% sans-serif
% \usepackage[sfdefault]{FiraSans}
% \usepackage[sfdefault]{roboto}
% \usepackage[sfdefault]{noto-sans}
% \usepackage[default]{sourcesanspro}

% serif
% \usepackage{CormorantGaramond}
% \usepackage{charter} % no T2A/Cyrillic support with pdfLaTeX
\usepackage{fontawesome5}


\pagestyle{fancy}
\fancyhf{} % clear all header and footer fields
\fancyfoot{}
\renewcommand{\headrulewidth}{0pt}
\renewcommand{\footrulewidth}{0pt}
\setlength{\footskip}{4.1pt}

% Adjust margins
\addtolength{\oddsidemargin}{-0.5in}
\addtolength{\evensidemargin}{-0.5in}
\addtolength{\textwidth}{1in}
\addtolength{\topmargin}{-.5in}
\addtolength{\textheight}{1.0in}

\urlstyle{same}

\raggedbottom
\raggedright
\setlength{\tabcolsep}{0in}

% Sections formatting
\titleformat{\section}{
  \vspace{-4pt}\scshape\raggedright\large
}{}{0em}{}[\color{black}\titlerule \vspace{-5pt}]

% Ensure that generate pdf is machine readable/ATS parsable
\ifPDFTeX
  \pdfgentounicode=1
\fi

%-------------------------
% Custom commands
\newcommand{\resumeItem}[1]{
  \item\small{
    {#1 \vspace{-2pt}}
  }
}

\newcommand{\resumeSubheading}[4]{
  \vspace{-2pt}\item
    \begin{tabular*}{0.97\textwidth}[t]{l@{\extracolsep{\fill}}r}
      \textbf{#1} & #2 \\
      \textit{\small#3} & \textit{\small #4} \\
    \end{tabular*}\vspace{-7pt}
}

\newcommand{\resumeSubSubheading}[2]{
    \item
    \begin{tabular*}{0.97\textwidth}{l@{\extracolsep{\fill}}r}
      \textit{\small#1} & \textit{\small #2} \\
    \end{tabular*}\vspace{-7pt}
}

\newcommand{\resumeProjectHeading}[2]{
    \item
    \begin{tabular*}{0.97\textwidth}{l@{\extracolsep{\fill}}r}
      \small#1 & #2 \\
    \end{tabular*}\vspace{-7pt}
}

\newcommand{\resumeSubItem}[1]{\resumeItem{#1}\vspace{-4pt}}

\renewcommand\labelitemii{$\vcenter{\hbox{\tiny$\bullet$}}$}

\newcommand{\resumeSubHeadingListStart}{\begin{itemize}[leftmargin=0.15in, label={}]}
\newcommand{\resumeSubHeadingListEnd}{\end{itemize}}
\newcommand{\resumeItemListStart}{\begin{itemize}}
\newcommand{\resumeItemListEnd}{\end{itemize}\vspace{-5pt}}

%-------------------------------------------
%%%%%%  RESUME STARTS HERE  %%%%%%%%%%%%%%%%%%%%%%%%%%%%


\begin{document}

%----------HEADING----------

\begin{center}

{\Huge \textbf{Станислав Макаров}}

\vspace{4pt}

{\large Senior QA Automation Engineer \textbar\ Python \textbar\ SDET}

\vspace{4pt}

Москва, Россия \textbullet\ Готов к релокации

\vspace{6pt}

\small
+7 993 250-39-60 \textbar\
\href{mailto:takahiroMushashii@proton.me}{takahiroMushashii@proton.me}
\textbar\
\href{https://github.com/chyioIshi}{GitHub}
\textbar\
\href{https://linkedin.com/in/chyioIshi}{LinkedIn}
\textbar\
\href{https://t.me/makarovsvv}{Telegram}

\end{center}

\section{О себе}
    \resumeSubHeadingListStart
        \item{\textbf{QA Automation Engineer / SDET} с \textbf{3+ годами опыта разработки автотестов и тестовых инструментов для микросервисных платформ}.
        Основная специализация -- \textbf{автоматизация на Python}, \textbf{API-тестирование}, \textbf{UI-тестирование} и \textbf{разработка тестовой инфраструктуры}.
        Работаю в \textbf{продуктовых командах}.

        Сейчас работаю в \textbf{Ozon}, занимаюсь разработкой автотестов и поддержкой тестовой инфраструктуры для внутренней платформы профиля сотрудника.
        Ранее в \textbf{Альфа-Банке} с нуля разработал \textbf{Python-экосистему тестирования, включающую API/UI-фреймворки, mock-сервис (FastAPI), CI/CD-пайплайн
        и модель автотестов (~1200 тестов)} для интеграционного тестирования микросервисной архитектуры.}
    \resumeSubHeadingListEnd

\section{Technical Stack}

\begin{itemize}[leftmargin=0.15in, label={}]
\small{\item{

\textbf{Programming:}
Python, SQL \\

\textbf{Test Automation:}
pytest, Playwright, Nuke, Allure TestOps, Jira Zephyr \\

\textbf{Backend:}
FastAPI, requests, Pydantic v2, Beanie ODM, PyMongo, kafka-python \\

\textbf{API \& Integration:}
REST API, gRPC, GraphQL, SOAP, Kafka, OpenAPI \\

\textbf{Databases:}
PostgreSQL, MongoDB, Valkey \\

\textbf{DevOps:}
Git, GitLab CI, Jenkins, Docker, Kubernetes, OpenShift, Helm, uv \\

\textbf{Observability:}
ELK Stack, Kibana, Prometheus, Grafana \\

\textbf{Architecture:}
Microservices, SOA, Event-Driven Architecture, Distributed Systems

}}
\end{itemize}

%-----------EXPERIENCE-----------
\section{Опыт работы}
  \resumeSubHeadingListStart

    \resumeSubheading
        {Ozon}{Май 2026 -- настоящее время}
        {Senior QA Automation Engineer}{Профиль сотрудника Ozon Job}

        \resumeItemListStart

            \resumeItem{\textbf{Проект:} Профиль сотрудника -- кластер backend сервисов, предоставляющих доступ для управления профилями сотрудников, финансовой информацией, сменами, данными о курсах и статусными моделями.}

            \resumeItem{Разрабатываю и поддерживаю end-to-end, API и интеграционные тесты с использованием Python, pytest и внутреннего фреймворка Nuke.}

            \resumeItem{Развиваю существующий фреймворк автоматизированного тестирования и улучшаю поддерживаемость тестового кода.}

            \resumeItem{Автоматизирую регрессионное тестирование новой функциональности.}

            \resumeItem{Проектирую и реализую интеграционные тесты для сервисов, преимущественно взаимодействующих через gRPC, а также REST API и Kafka.}

            \resumeItem{Исследую дефекты и проблемы на тестовых окружениях, провожу первичный анализ причин и взаимодействовую с разработчиками для их устранения.}

            \resumeItem{Работаю с \textbf{Kubernetes}, \textbf{GitLab CI}, \textbf{Valkey}, \textbf{шардированным PostgreSQL} при проверке поведения распределенной системы.}

            \resumeItem{Участвую в проверке функциональности на разных этапах разработки, поддерживая стабильную поставку новых изменений.}

        \resumeItemListEnd

    \resumeSubheading
        {Альфа-Банк}{Март 2025 -- Май 2026}
        {Senior QA Automation Engineer / Python Tooling Developer}
        {Валютные деривативы НИБ}

        \resumeItemListStart

            \resumeItem{\textbf{Проект:} кластер UI/API сервисов, взаимодействующих через REST API, Kafka, внутренние сервисы и
            осуществляющих поддержку создания и проведения и получения информации о сделках (расчётные и процентные форварды, свопы, споты) для ЮЛ.}

            \resumeItem{Разработал mock-сервис на FastAPI для эмуляции внешних систем, что позволило проводить изолированное интеграционное тестирование и гибко покрывать негативные сценарии.}

            \resumeItem{Создал внутренний Python-сервис (Zeep, SOAP, REST API, MongoDB) для автоматизированной подготовки тестовых пользователей через ESB, сократив время подготовки тестовых данных с \textbf{1--2 дней примерно до 10 минут}.}

            \resumeItem{Построил \textbf{CI/CD-пайплайн на Groovy (Jenkins)} для запуска автотестов, публикации отчётов Allure и интеграции тестирования в процесс поставки релизов.}

            \resumeItem{С нуля разработал \textbf{фреймворк автоматизации тестирования} для API, интеграционного и UI-тестирования микросервисной архитектуры.}

            \resumeItem{Автоматизировал \textbf{до 90\% тестов новой функциональности}, включая \textbf{интеграционные, контрактные и UI-тесты}.}

            \resumeItem{Разработал стандарт структурированного логирования сервисов (JSON logs, trace-id, request/response metadata), что ускорило локализацию дефектов.}

            \resumeItem{Выполнял диагностику проблем системы с использованием \textbf{ELK, Kibana, Prometheus и Grafana}.}

            \resumeItem{Работал с \textbf{Kubernetes} (pods, services, ingresses, deployments, network policies, ansible-скрипты, Helm Charts) при анализе микросервисов на 
            тестовых стендах.}

            \resumeItem{Активно участвовал в \textbf{Shift Left Testing}, анализируя требования и архитектуру ещё до начала разработки.}

        \resumeItemListEnd

    \resumeSubheading
        {Сбер}{Октябрь 2023 -- Март 2025}
        {QA Automation Engineer}
        {Виртуальный ассистент}

        \resumeItemListStart

            \resumeItem{\textbf{Проект:} Виртуальный голосовой помощник «Салют» в сфере инвестиций и накоплений - позволяет клиенту проводить активные операции
            голосом, получать персональные консультации по сложным инвестиционнонакопительным продуктам с использованием AI-моделей и интеграцией с GigaChat.}

            \resumeItem{Увеличил тестовое покрытие системы \textbf{с 65\% до 80\%}, разработав дополнительные \textbf{модульные и интеграционные тесты}.}

            \resumeItem{Провёл \textbf{рефакторинг тестового фреймворка}, внедрив паттерны проектирования и повысив поддерживаемость тестового кода.}

            \resumeItem{Автоматизировал тестирование \textbf{AI-модели классификации пользовательских намерений}, обнаружив критический дефект между моделями DEV и PROD.}

            \resumeItem{Повысил уровень автоматизации \textbf{с 60\% до 90\%}, существенно ускорив прохождение регрессионного тестирования.}

            \resumeItem{Разработал \textbf{интеграционные тесты для микросервисов}, взаимодействующих через \textbf{Kafka и REST API}.}

            \resumeItem{Менторил junior QA и стажёров в области \textbf{автоматизации тестирования и CI/CD-процессов}.}

        \resumeItemListEnd

  \resumeSubHeadingListEnd


%-----------PROJECTS-----------
\section{Проекты}

\resumeSubHeadingListStart

\resumeProjectHeading
{\textbf{DevMirror}
\href{https://github.com/chyioIshi/DevMirror}{\raisebox{0.2ex}{\footnotesize\faGithub}}
$|$
\emph{Python, FastAPI, MongoDB, Kafka, Docker}}
{}

\resumeItemListStart

\resumeItem{Спроектировал и реализовал универсальный mock-сервис для интеграционного тестирования распределенных микросервисных систем.}

\resumeItem{Разработал сервис с применением принципов DDD и Clean Architecture.}

\resumeItem{Реализовал динамическое сопоставление запросов, scoped mocks, side effects, логирование операций и интеграцию с Kafka.}

\resumeItem{Разработал расширяемый side-effect engine с поддержкой HTTP callbacks, PostgreSQL, MongoDB, Redis и Kafka providers.}

\resumeItem{Реализовал CI-пайплайн в GitHub Actions с ruff/mypy linting и unit/integration тестами.}

\resumeItemListEnd

\resumeSubHeadingListEnd



\section{Образование}

    \resumeSubHeadingListStart

    \resumeSubheading
    {МИРЭА -- Российский технологический университет}
    {Москва, Россия}
    {Бакалавриат, лазерная техника и лазерные технологии}
    {}

    \resumeSubHeadingListEnd


%-------------------------------------------
\end{document}
